% Источник доски: принцип «переводим в степенной вид:
% составные числа, корни, дроби». :contentReference[oaicite:0]{index=0}

\documentclass[a4paper,12pt]{article}

% ------------------------------------------------
% Кодировка, язык, математика
% ------------------------------------------------
\usepackage[T2A]{fontenc}
\usepackage[utf8]{inputenc}
\usepackage[russian]{babel}

\usepackage{amsmath,amssymb,amsthm,mathtools}
\usepackage{icomma}

% ------------------------------------------------
% Типографика
% ------------------------------------------------
\usepackage{helvet}
\renewcommand{\familydefault}{\sfdefault}
\usepackage{microtype}
\usepackage{setspace}
\onehalfspacing
\usepackage{parskip}

% ------------------------------------------------
% Страница
% ------------------------------------------------
\usepackage[
  a4paper,
  left=20mm,
  right=20mm,
  top=20mm,
  bottom=22mm
]{geometry}

% ------------------------------------------------
% Графика
% ------------------------------------------------
\usepackage{tikz}
\usetikzlibrary{calc}
\usetikzlibrary{arrows.meta,positioning,shapes.geometric}

% ------------------------------------------------
% Таблицы и списки
% ------------------------------------------------
\usepackage{tabularx,booktabs,longtable}
\usepackage{enumitem}
\renewcommand{\arraystretch}{1.35}

% ------------------------------------------------
% Колонтитулы и ссылки
% ------------------------------------------------
\usepackage{fancyhdr}
\usepackage{hyperref}
\hypersetup{
  hidelinks,
  unicode=true,
  pdftitle={Степени в прикладных задачах},
  pdfauthor={Лёвин Артём Александрович}
}

% ------------------------------------------------
% Цвета
% ------------------------------------------------
\usepackage{xcolor}
\definecolor{accent}{HTML}{008080}
\definecolor{darkaccent}{HTML}{004D4D}
\definecolor{textgray}{HTML}{333333}
\definecolor{softgray}{HTML}{E8EEEE}

\color{textgray}

% ------------------------------------------------
% Колонтитулы
% ------------------------------------------------
\pagestyle{fancy}
\fancyhf{}
\fancyhead[L]{\small\color{darkaccent}Степени в прикладных задачах}
\fancyhead[R]{\small\color{textgray}ЕГЭ}
\fancyfoot[C]{\small\thepage}
\renewcommand{\headrulewidth}{0.3pt}
\renewcommand{\headrule}{%
  \hbox to\headwidth{\color{accent}\leaders\hrule height \headrulewidth\hfill}
}

% ------------------------------------------------
% Списки
% ------------------------------------------------
\setlist[itemize]{
  leftmargin=6mm,
  itemsep=4pt,
  topsep=4pt
}

\setlist[enumerate]{
  leftmargin=8mm,
  itemsep=6pt,
  topsep=5pt
}

% ------------------------------------------------
% Служебные команды
% ------------------------------------------------
\newcommand{\topic}[1]{%
  \vspace{0.5em}
  {\Large\bfseries\color{darkaccent}#1}\par
  \vspace{0.25em}
  {\color{accent}\rule{\linewidth}{0.7pt}}\par
  \vspace{0.6em}
}

\newcommand{\subtopic}[1]{%
  \vspace{0.5em}
  {\large\bfseries\color{darkaccent}#1}\par
  \vspace{0.25em}
}

\newcommand{\keyidea}[1]{%
  \par\medskip
  {\color{accent}\rule{8mm}{1.2pt}}\quad
  \textbf{#1}
  \par\smallskip
}

% ------------------------------------------------
% Стили блок-схемы
% ------------------------------------------------
\tikzset{
  flowstart/.style={
    ellipse,
    draw=accent,
    line width=0.9pt,
    minimum width=46mm,
    minimum height=10mm,
    align=center,
    inner sep=3pt,
    font=\small
  },
  flowstep/.style={
    rectangle,
    rounded corners=3pt,
    draw=accent,
    line width=0.8pt,
    text width=125mm,
    minimum height=11mm,
    align=center,
    inner sep=4pt,
    font=\small
  },
  flowarrow/.style={
    -{Latex[length=2.5mm,width=1.7mm]},
    line width=0.8pt,
    draw=darkaccent
  }
}

\begin{document}

% ================================================================
% ТИТУЛЬНЫЙ ЛИСТ
% ================================================================
\begin{titlepage}
\thispagestyle{empty}

\begin{center}

\vspace*{22mm}

{\small\color{accent}\bfseries
ПРОФИЛЬНАЯ МАТЕМАТИКА \(\cdot\) ЕГЭ
}

\vspace{18mm}

{\Huge\bfseries\color{darkaccent}
Чек-лист
}

\vspace{7mm}

{\LARGE\bfseries
Степени в прикладных задачах
}

\vspace{8mm}

{\large
Работа с формулами, степенной записью,\\
большими числами и функциональными зависимостями
}

\vfill

{\large\bfseries
Екатерина Гнедкова
}

\vspace{6mm}

{\normalsize
\today
}

\vspace{13mm}

{\small
Лёвин Артём Александрович\\
эксклюзивно для Екатерины Гнедковой
}

\end{center}

\begin{tikzpicture}[remember picture,overlay]
  \draw[
    accent,
    line width=1.1pt
  ]
  ($(current page.south west)+(16mm,16mm)$)
  .. controls
  ($(current page.south west)+(62mm,28mm)$)
  and
  ($(current page.south east)+(-62mm,8mm)$)
  ..
  ($(current page.south east)+(-16mm,17mm)$);
\end{tikzpicture}

\end{titlepage}

% ================================================================
% ОСНОВНОЙ ЧЕК-ЛИСТ
% ================================================================
\topic{1. Главная идея занятия}

Прикладная задача часто выглядит объёмной из-за большого текста,
обозначений и чисел. После перевода условия на язык формул основная
математика обычно сводится к работе со степенями.

\keyidea{Главный принцип}

Если в ходе решения появляется степень, стремитесь представить
остальные элементы выражения в удобном степенном виде.

Особенно важно преобразовывать:

\begin{enumerate}
  \item составные числа;
  \item корни;
  \item дроби.
\end{enumerate}

Например,
\[
216=2^3\cdot3^3,
\qquad
1000=10^3,
\qquad
\frac18=2^{-3},
\qquad
625=5^4.
\]

Такое представление позволяет сокращать множители, складывать и
вычитать показатели степеней и быстро извлекать требуемую степень.

% ================================================================
\topic{2. Базовые свойства степеней}

Для допустимых значений переменных:

\[
a^m\cdot a^n=a^{m+n},
\]

\[
\frac{a^m}{a^n}=a^{m-n},
\qquad a\ne0,
\]

\[
\left(a^m\right)^n=a^{mn},
\]

\[
(ab)^n=a^n b^n,
\]

\[
\left(\frac ab\right)^n=\frac{a^n}{b^n},
\qquad b\ne0.
\]

Отрицательная степень:

\[
a^{-n}=\frac1{a^n},
\qquad a\ne0.
\]

Дробная степень при положительном основании:

\[
a^{\frac mn}
=
\sqrt[n]{a^m}.
\]

Если получено

\[
x^r=A,
\]

для выделения \(x\) удобно возвести обе части в степень,
обратную \(r\):

\[
x=\left(A\right)^{1/r}.
\]

Для физических величин дополнительно учитывается их смысл:
длина, объём, масса, температура и подобные величины обычно
рассматриваются как положительные.

% ================================================================
\topic{3. Как работать с данными прикладной задачи}

\subtopic{Шаг 1. Выпишите рабочую формулу}

Сначала отделите формулу от большого текстового условия.

Пример структуры:

\[
P=\sigma ST^4.
\]

После этого отдельно выпишите известные величины и искомую величину.

\subtopic{Шаг 2. Приведите числа к удобному виду}

При вычислениях со степенями десятичную запись часто выгодно
переписать так, чтобы коэффициент стал целым:

\[
5{,}7\cdot10^{-8}
=
57\cdot10^{-9},
\]

\[
1{,}14\cdot10^{25}
=
114\cdot10^{23}.
\]

Правило:

\begin{itemize}
  \item запятая сдвинулась вправо на один разряд
  \(\Rightarrow\) показатель десятки уменьшился на \(1\);
  \item запятая сдвинулась вправо на два разряда
  \(\Rightarrow\) показатель десятки уменьшился на \(2\).
\end{itemize}

Проверка на первом примере:

\[
5{,}7=\frac{57}{10},
\]

поэтому

\[
5{,}7\cdot10^{-8}
=
57\cdot10^{-1}\cdot10^{-8}
=
57\cdot10^{-9}.
\]

\subtopic{Шаг 3. Следите за единицами}

Внутри длинных вычислений единицы измерения можно не повторять
в каждой строке. До подстановки все величины должны быть приведены
к единицам, которые предполагает формула.

Например, если формула требует метры, значение в миллиметрах
сначала переводится в метры.

% ================================================================
\topic{4. Слова «не больше», «не меньше», «наибольший»}

В прикладных задачах такие формулировки часто задают граничное
значение допустимого диапазона.

Если требуется наибольшее или наименьшее значение и зависимость
монотонна на рассматриваемой области, искомая граница достигается
при равенстве.

Например, условие

\[
F\le1000
\]

при поиске максимальной длины \(L\) может привести к уравнению

\[
F=1000.
\]

Перед переходом к равенству следует убедиться, что:

\begin{itemize}
  \item рассматриваемые величины лежат в допустимой области;
  \item формула действительно задаёт нужную монотонность;
  \item вопрос задачи относится к граничному значению.
\end{itemize}

% ================================================================
\topic{5. Пример: кубическая зависимость}

После подстановки чисел получено уравнение

\[
1000=10000L^3-250.
\]

Переносим число:

\[
1250=10000L^3.
\]

Выражаем куб:

\[
L^3=\frac{1250}{10000}
=\frac{125}{1000}
=\frac18.
\]

Переводим правую часть в степенной вид:

\[
\frac18=\left(\frac12\right)^3.
\]

Следовательно,

\[
L=\frac12=0{,}5.
\]

\keyidea{Что здесь важно}

Как только возникло \(L^3\), правая часть была представлена как
куб удобного числа.

% ================================================================
\topic{6. Пример: четвёртая степень и большие числа}

Пусть

\[
P=\sigma ST^4,
\]

где

\[
\sigma=5{,}7\cdot10^{-8}
      =57\cdot10^{-9},
\]

\[
S=\frac{10^{20}}{128},
\]

\[
P=1{,}14\cdot10^{25}
 =114\cdot10^{23}.
\]

Подставляем данные:

\[
114\cdot10^{23}
=
57\cdot10^{-9}
\cdot
\frac{10^{20}}{128}
\cdot T^4.
\]

Выражаем \(T^4\):

\[
T^4
=
\frac{
114\cdot10^{23}\cdot128
}{
57\cdot10^{-9}\cdot10^{20}
}.
\]

Работаем отдельно с числовыми множителями и степенями десятки:

\[
\frac{114\cdot128}{57}
=
2\cdot128
=
256
=
2^8,
\]

\[
10^{23}\cdot10^9\cdot10^{-20}
=
10^{12}.
\]

Получаем

\[
T^4=2^8\cdot10^{12}.
\]

Замечаем:

\[
2^8\cdot10^{12}
=
\left(2^2\cdot10^3\right)^4.
\]

Так как температура положительна,

\[
T=2^2\cdot10^3=4000.
\]

\keyidea{Почему факторизация полезна}

Числа \(114\), \(57\), \(128\) выглядят громоздко.
После разложения и сокращения вычисление превращается в работу
с показателями степеней.

% ================================================================
\topic{7. Пример: функциональная зависимость}

Пусть масса изменяется по закону

\[
m(t)=m_0\cdot2^{-t/T}.
\]

Запись \(m(t)\) означает, что масса является функцией времени \(t\).

Пусть

\[
m_0=40,
\qquad
T=10,
\qquad
m(t)=5.
\]

Подставляем:

\[
5=40\cdot2^{-t/10}.
\]

Делим обе части на \(40\):

\[
\frac18=2^{-t/10}.
\]

Переводим дробь в степенной вид:

\[
\frac18=2^{-3}.
\]

При одинаковых положительных основаниях показатели равны:

\[
-\frac{t}{10}=-3.
\]

Следовательно,

\[
t=30.
\]

\keyidea{Полезная привычка}

Дроби вида

\[
\frac12,\quad
\frac14,\quad
\frac18,\quad
\frac1{16}
\]

сразу распознавайте как

\[
2^{-1},\quad
2^{-2},\quad
2^{-3},\quad
2^{-4}.
\]

% ================================================================
\topic{8. Пример: закон с постоянной величиной}

Пусть

\[
pV^a=\mathrm{const}.
\]

Эта запись означает, что при переходе от одного состояния к другому

\[
p_1V_1^a=p_2V_2^a.
\]

Предположим, что объём уменьшился в \(2\) раза:

\[
V_2=\frac{V_1}{2}.
\]

Для подстановки удобно переписать:

\[
V_1=2V_2.
\]

Давление увеличилось в \(4\) раза:

\[
p_2=4p_1.
\]

Подставляем:

\[
p_1(2V_2)^a
=
4p_1V_2^a.
\]

Раскрываем степень произведения:

\[
p_1\cdot2^a\cdot V_2^a
=
4p_1V_2^a.
\]

При положительных \(p_1\) и \(V_2\) сокращаем одинаковые множители:

\[
2^a=4.
\]

Так как

\[
4=2^2,
\]

получаем

\[
\boxed{a=2}.
\]

\keyidea{Перевод текста на язык математики}

Фразы

\[
\text{«уменьшился в 2 раза»}
\]

и

\[
\text{«увеличился в 4 раза»}
\]

следует сразу превращать в отношения между старым и новым
значением величины.

% ================================================================
% БЛОК-СХЕМА
% ================================================================
\clearpage
\thispagestyle{fancy}

\begin{center}
{\LARGE\bfseries\color{darkaccent}
Алгоритм: прикладная задача со степенями
}
\end{center}

\vspace{6mm}

\begin{center}
\begin{tikzpicture}[node distance=7mm]

\node[flowstart] (start)
{Начало};

\node[flowstep,below=of start] (read)
{Прочитайте условие, выделите рабочую формулу,
известные величины и искомую величину};

\node[flowstep,below=of read] (prepare)
{Приведите единицы измерения и числовую запись
к удобному виду};

\node[flowstep,below=of prepare] (boundary)
{Если задача просит граничное значение,
обоснуйте переход от ограничения к равенству};

\node[flowstep,below=of boundary] (substitute)
{Подставьте данные и выразите степень искомой величины:
\(x^r=A\)};

\node[flowstep,below=of substitute] (powerform)
{Переведите составные числа, корни и дроби
в удобный степенной вид; разложите крупные числа на множители};

\node[flowstep,below=of powerform] (simplify)
{Сократите одинаковые множители и примените
свойства степеней};

\node[flowstep,below=of simplify] (extract)
{Возведите обе части в степень \(1/r\)
и получите искомую величину};

\node[flowstep,below=of extract] (check)
{Проверьте область допустимых значений,
знак ответа и соответствие смыслу задачи};

\node[flowstart,below=of check] (finish)
{Ответ};

\draw[flowarrow] (start) -- (read);
\draw[flowarrow] (read) -- (prepare);
\draw[flowarrow] (prepare) -- (boundary);
\draw[flowarrow] (boundary) -- (substitute);
\draw[flowarrow] (substitute) -- (powerform);
\draw[flowarrow] (powerform) -- (simplify);
\draw[flowarrow] (simplify) -- (extract);
\draw[flowarrow] (extract) -- (check);
\draw[flowarrow] (check) -- (finish);

\end{tikzpicture}
\end{center}

\clearpage

% ================================================================
\topic{9. Типичные ошибки}

\begin{enumerate}

\item
\textbf{Неверное изменение порядка десятки.}

Проверяйте:

\[
5{,}7\cdot10^{-8}
=
57\cdot10^{-9}.
\]

Коэффициент увеличился в \(10\) раз, поэтому степень десятки
уменьшилась на единицу.

\item
\textbf{Попытка считать большие числа напрямую.}

При степенной задаче сначала ищите факторизацию:

\[
128=2^7,
\qquad
625=5^4,
\qquad
216=2^3\cdot3^3.
\]

\item
\textbf{Потеря обратной степени.}

Из

\[
x^{5/3}=A
\]

следует

\[
x=A^{3/5}.
\]

\item
\textbf{Неверное чтение записи функции.}

\[
m(t)
\]

обозначает значение функции \(m\) в момент \(t\).

\item
\textbf{Автоматический переход от неравенства к равенству.}

Для граничной задачи такой переход должен соответствовать
монотонности зависимости и смыслу условия.

\item
\textbf{Игнорирование допустимых значений.}

Например, из

\[
x^4=16
\]

в обычном алгебраическом уравнении следуют два корня:

\[
x=\pm2.
\]

Если \(x\) обозначает положительную физическую величину,
остаётся

\[
x=2.
\]

\end{enumerate}

% ================================================================
\topic{10. Короткий чек-лист перед ответом}

Перед завершением задачи проверьте:

\begin{itemize}
  \item формула выписана верно;
  \item все необходимые данные найдены в тексте;
  \item единицы согласованы;
  \item десятичные множители приведены к удобной записи;
  \item составные числа, корни и дроби переведены в степенной вид;
  \item показатели степеней обработаны по правилам;
  \item искомая величина выделена;
  \item учтён её допустимый знак;
  \item ответ соответствует смыслу исходной задачи.
\end{itemize}

% ================================================================
% ТРЕНИРОВОЧНЫЙ БЛОК
% ================================================================
\clearpage

\begin{center}
{\LARGE\bfseries\color{darkaccent}
Тренировочный блок
}

\vspace{2mm}

{\large
5 заданий \(\cdot\) от простого к сложному
}
\end{center}

\vspace{5mm}

\begin{enumerate}

\item
После подстановки данных в формулу для некоторой положительной
длины \(L\) получено уравнение

\[
840=8000L^3-160.
\]

Найдите \(L\).

\item
Масса вещества изменяется по закону

\[
m(t)=64\cdot2^{-t/6}.
\]

Через сколько единиц времени масса станет равной \(8\)?

\item
Для положительных величин \(p\) и \(V\) выполняется

\[
pV^{3/2}=27.
\]

Найдите \(V\), если

\[
p=8.
\]

\item
Излучаемая мощность определяется формулой

\[
P=\sigma ST^4.
\]

Известно:

\[
\sigma=3{,}2\cdot10^{-8},
\qquad
S=5\cdot10^{16},
\qquad
P=10^{24}.
\]

Найдите положительное значение \(T\).

\item
Для некоторого процесса выполняется закон

\[
pV^a=\mathrm{const}.
\]

При переходе из первого состояния во второе объём уменьшился
в \(3\) раза, а давление увеличилось в \(9\) раз.

Найдите показатель \(a\).

\end{enumerate}

% ================================================================
% ОТВЕТЫ
% ================================================================
\clearpage

\begin{center}
{\LARGE\bfseries\color{darkaccent}
Ответы к тренировочному блоку
}
\end{center}

\vspace{8mm}

\begin{table}[ht]
\centering
\begin{tabularx}{\linewidth}{@{}>{\centering\arraybackslash}p{24mm}X@{}}
\toprule
\textbf{№} & \textbf{Ответ} \\
\midrule
1 & \(L=\dfrac12=0{,}5\) \\[2mm]
2 & \(t=18\) \\[2mm]
3 & \(V=\dfrac94=2{,}25\) \\[2mm]
4 & \(T=5000\) \\[2mm]
5 & \(a=2\) \\
\bottomrule
\end{tabularx}
\end{table}

\vfill

\begin{center}
{\small\color{textgray}
Лёвин Артём Александрович\\
эксклюзивно для Екатерины Гнедковой
}
\end{center}

\end{document}